%==============================================================================
% Scientific manuscript template — self-contained, journal-neutral.
% Compiles with: latexmk -pdf -interaction=nonstopmode paper.tex
% Requires only standard TeX Live packages (no journal class, no biblatex).
%
% To target a specific journal, swap \documentclass and add their metadata
% macros; the body should transfer unchanged.
%==============================================================================
\documentclass[11pt,a4paper]{article}

% --- Encoding & language -----------------------------------------------------
\usepackage[T1]{fontenc}
\usepackage[utf8]{inputenc}

% --- Fonts -------------------------------------------------------------------
\usepackage{mathptmx}          % Times text + matching math. Alternative: charter
\usepackage{microtype}         % Fixes most justification complaints automatically

% --- Page geometry -----------------------------------------------------------
\usepackage[margin=25mm]{geometry}

% --- Mathematics -------------------------------------------------------------
\usepackage{amsmath,amssymb,mathtools}
\usepackage{bm}
\usepackage{siunitx}           % \SI{3.2}{\micro\meter}, S columns for tables

% --- Floats ------------------------------------------------------------------
\usepackage{graphicx}
\usepackage{booktabs}          % \toprule \midrule \bottomrule — no vertical rules
\usepackage{threeparttable}    % table footnotes that stay inside the float
\usepackage{multirow}
\usepackage{longtable}
\usepackage{float}
\usepackage{placeins}          % \FloatBarrier
\usepackage[font=small,labelfont=bf,width=0.92\textwidth]{caption}
\usepackage{subcaption}

% --- Structure ---------------------------------------------------------------
\usepackage{titlesec}
\usepackage{enumitem}
\usepackage{setspace}          % \onehalfspacing for review copies
\usepackage{lineno}            % \linenumbers when a journal wants them

% --- References & links (load order matters: hyperref late, cleveref last) ----
\usepackage[numbers,sort&compress]{natbib}
\bibliographystyle{plainnat}   % or unsrtnat, abbrvnat, apalike, ieeetr
\usepackage{url}
\usepackage[hidelinks]{hyperref}
\usepackage[capitalise,nameinlink]{cleveref}

% --- Typographic hygiene -----------------------------------------------------
\widowpenalty=10000            % no single line stranded at a page top
\clubpenalty=10000             % no single line stranded at a page bottom
\displaywidowpenalty=10000
\sloppy\emergencystretch=1em   % prefer slightly loose lines over margin overflow

\titleformat{\section}{\normalfont\fontsize{14}{16}\bfseries}{\thesection}{1em}{}
\titleformat{\subsection}{\normalfont\fontsize{12}{14}\bfseries}{\thesubsection}{1em}{}
\titleformat{\subsubsection}{\normalfont\fontsize{11}{13}\itshape}{\thesubsubsection}{1em}{}
\titlespacing*{\section}{0pt}{14pt plus 3pt minus 2pt}{6pt plus 2pt minus 1pt}
\titlespacing*{\subsection}{0pt}{11pt plus 2pt minus 2pt}{4pt plus 2pt minus 1pt}

%==============================================================================
\begin{document}

\title{\bfseries Full Title Stating the Finding, Not the Topic}
\author{Author Name\thanks{Affiliation. Corresponding author: email@example.org}}
\date{\today}
\maketitle

\begin{abstract}
\noindent
Context in one or two sentences. The specific question. What was done. What was
found, including the key quantity with its uncertainty. What it means. The
principal limitation. Written last, from \texttt{claims.json}; only
source-backed or explicitly hedged results appear here.
\end{abstract}

\vspace{0.5em}
\noindent\textbf{Keywords:} keyword one; keyword two; keyword three

%------------------------------------------------------------------------------
\section{Introduction}
\label{sec:intro}

Four moves: the problem, what is known, the gap, the contribution. The gap
follows from the search; the contribution is sized to the evidence
\citep{example2024}.

%------------------------------------------------------------------------------
\section{Background}
\label{sec:background}

Formalism and prior results the reader needs. Established results are cited and
labelled as known:
\begin{equation}
  S = k_B \ln \Omega ,
  \label{eq:boltzmann}
\end{equation}
where $S$ is the entropy (\si{\joule\per\kelvin}), $k_B$ the Boltzmann constant,
and $\Omega$ the number of accessible microstates. \Cref{eq:boltzmann} is a
standard result \citep{example2024}; equations derived in this work are
identified as such where they appear.

%------------------------------------------------------------------------------
\section{Methods}
\label{sec:methods}

Enough detail to reproduce: procedure, parameters, seeds, software versions,
and the analysis choices that were fixed in advance.

%------------------------------------------------------------------------------
\section{Results}
\label{sec:results}

Findings only. Interpretation belongs in the discussion.

\begin{figure}[!htbp]
  \centering
  % \includegraphics[width=0.8\linewidth]{figures/fig1.pdf}
  \rule{0.8\linewidth}{4cm}% placeholder — replace with \includegraphics
  \caption{Self-contained caption. State what is plotted, the units, the
           uncertainty convention, and $n$. A reader who reads only the figures
           should understand this one.}
  \label{fig:example}
\end{figure}

\Cref{fig:example} shows the primary result, and \cref{tab:example} reports the
fitted values. Every float must be referenced in the text like this --
\texttt{pdf\_qa.py} treats an undiscussed figure or table as an error, because it
means either the float is unnecessary or the text has a gap.

\begin{table}[!htbp]
  \centering
  \caption{Fit results across binning choices.}
  \label{tab:example}
  \begin{tabular}{l S[table-format=1.2] S[table-format=1.2]}
    \toprule
    Bins & {Exponent} & {Uncertainty} \\
    \midrule
    20  & 1.83 & 0.07 \\
    50  & 1.84 & 0.06 \\
    100 & 1.85 & 0.08 \\
    \bottomrule
  \end{tabular}
\end{table}

\FloatBarrier

%------------------------------------------------------------------------------
\section{Discussion}
\label{sec:discussion}

Interpretation, with inferences marked as inferences and alternative
explanations considered rather than dismissed.

%------------------------------------------------------------------------------
\section{Limitations}
\label{sec:limitations}

Specific, not ritual. State what would change the conclusion.

%------------------------------------------------------------------------------
\section{Conclusion}
\label{sec:conclusion}

What is now known that was not known before, at the size the evidence supports.

%------------------------------------------------------------------------------
\section*{Data and code availability}
Where the data and analysis code live, and under what terms.

\section*{Disclosure of AI involvement}
This manuscript was prepared with substantial assistance from an AI system
(literature retrieval, drafting, analysis, and typesetting). All claims were
checked against retrieved sources, and all reported computations were executed.
Responsibility for the content rests with the named author(s).

\section*{Acknowledgements}
Funding, facilities, and people.

%------------------------------------------------------------------------------
\bibliography{refs}

\appendix
\section{Supplementary derivations}
\label{app:derivations}

Routine algebra and secondary analyses.

\end{document}
